\scanpage{287}
\mspage{237}

déduites d'inclusions $A\subset A'$ \unsure{dans|ds} $\unsure{P|\mathcal P}(Ob\,I)$. Ainsi, le caractère $W$-cocartésien des carrés 1, 2, 3 s'\unsure{exprime|expriment} respectivement par la condition que les inclusions
\[
 \begin{aligned}
 \unsure{\mathscr X|\mathcal X}_{010'}&\xrightarrow{\varphi_1}\mathscr X_{010'1'}\\
 \mathscr X_{121'}&\xrightarrow{\varphi_2}\mathscr X_{121'2'}\\
 \mathscr X_{020'}&\xrightarrow{\varphi_3}\mathscr X_{020'2'}
 \end{aligned}
\]
soient dans $W$. Considérons \unsure{donc|alors} le diagramme de flèches d'inclusion
\diagram{\begin{tikzcd}[ampersand replacement=\&]
 \mathscr X_{020'} \arrow[r,"{\varphi_{\deleted{\unsure{2}}3}}"] \arrow[d,"{\alpha}"'] \& \mathscr X_{020'2'} \arrow[d,"{\beta}"] \\
 \mathscr X_U \arrow[r,"{\psi_3}"] \arrow[d,"{\deleted{\unsure{\varphi}}\psi_1}"'] \& \mathscr X \arrow[d,"{id}"] \\
 \mathscr X_V \arrow[r,"{\psi_2}"] \& \mathscr X\\
 \mathscr X_{\deleted{\unsure{020'}}121'} \arrow[r,"{\varphi_2}"] \arrow[u,"{\alpha'}"] \& \mathscr X_{\deleted{\unsure{020'2'}}121'2'} \arrow[u,"{\beta'}"']
\end{tikzcd}}
Je dis que dans ce diagramme on a
\[
 \alpha,\beta,\alpha',\beta'\ \in W \tag*{$(*)$}
\]
et
\[
 \varphi_1\in W\ \text{(i.e. carré 1 $W$-cocartésien)}
 \Longleftrightarrow\deleted{\unsure{\varphi}}\psi_1\in W. \tag*{$(**)$}
\]
Les relations $(*)$ impliquent déjà que
\[
 \begin{aligned}
 \varphi_2\in W&\Longleftrightarrow\psi_2\in W\\
 \varphi_3\in W&\Longleftrightarrow\psi_3\in W
 \end{aligned}
\]
et sous l'hyp. $W$-cocartésienne $(**)$, impliquant $\deleted{\unsure{\varphi}}\psi_1\in W$,

% archive sheet number: 287

